\chapter{One Backend for Every Model}
\label{chap:s2-9}

\section{Goal}

The model roster (fifteen models, approved 2026-11-12) is only a comparison if every model is measured
by the same instrument. Until now the harness had one backend, written for Qwen2.5-VL. This step splits
it in two: a shared backend that holds every piece of measuring logic once, and a small adapter per model
family that holds only what genuinely differs. The test of the step is that the shared backend, given the
Qwen2.5-VL adapter, reproduces the frozen backend exactly.

\section{What Was Done}

\textbf{The split.} \texttt{backends/hf.py} (\texttt{HFVLMBackend}) holds generation and timing, reply
parsing, black masking, token ablation and the attention read-out. \texttt{backends/families.py} holds
one adapter per family, with five duties:
\begin{itemize}
\item loading the model and its processor;
\item turning one image and one prompt into model inputs;
\item stating the frame a pixel box is given in (for models that answer in pixels);
\item locating the image tokens and the grid they form;
\item merging vision features into the token embeddings, for token ablation.
\end{itemize}
A family whose image tokens do not form one grid over the whole image (tiled encoders, or Llama~3.2's
cross-attention) reports no grid. Attention and token ablation are then marked not applicable for that
model rather than computed on a guessed layout, as the roster design requires (D4).

\textbf{The logic was moved, not rewritten.} The token-cell test and the attention average were lifted out
of the frozen backend into two functions, \texttt{cells\_to\_keep} and \texttt{aggregate\_attention}. CPU
tests check both against verbatim copies of the frozen loops: 200 random regions on three grid shapes for
the first, synthetic attention tensors for the second (exact equality).

\textbf{Coordinate conventions.} \texttt{frames.parse\_in\_convention} reads a reply's REGION in one of
four units and returns it in original-image pixels:
\begin{itemize}
\item \texttt{input\_px}: pixels of the resized input (Qwen2.5-VL); exactly the corrected path of
Chapter~\ref{chap:s2-8};
\item \texttt{norm1000}: 0--1000 over the image (Qwen2-VL, Qwen3-VL, InternVL, GLM);
\item \texttt{norm1}: 0--1 over the image (the LLaVA family);
\item \texttt{none}: verdict-only models; the region is dropped.
\end{itemize}
Each has a round-trip test on one known box. The prompt follows the convention (roster D2): only the
REGION line changes, and a test checks that the default prompt is byte-identical to the frozen one for
every rule.

\textbf{Behind a flag.} \texttt{XAI\_BACKEND} selects the class for the two frozen tiers. The default,
\texttt{legacy}, keeps the frozen \texttt{QwenBackend}, so nothing already reported can move;
\texttt{hf} routes them through the shared backend. New tiers have no frozen results and always take the
shared backend, and they are always read in their own convention whatever \texttt{XAI\_REGION\_FRAME} says.
Every model stage now gets its backend from one factory (\texttt{backends/load.py}), and its fingerprint
records the class and convention whenever it is not the frozen path.

\textbf{Registry.} Each tier now names its family, its box convention and its environment (A, the frozen
\texttt{transformers} 4.49; B, the current release). Two Qwen2.5-VL size arms are registered and pinned:
\texttt{qwen3b-bf16} and \texttt{qwen32b-4bit}. Llama~3.2-11B-Vision is downloaded (Meta approved access on
2026-11-13) and pinned; it gets its adapter at its own Stage~0.

\section{Results}

\textbf{The check.} \texttt{scripts/20\_backend\_parity.py} runs the same rows through both backends,
one after the other, in one process on one GPU. Rows are drawn at random (seed 0) from those with a stated
region, without regard to their score, as Chapter~\ref{chap:s2-8} requires. For each row it records:
\begin{itemize}
\item the reply to the clean image, its verdict and its region;
\item the reply with the region blacked out (Stage~4);
\item the reply with the region's image tokens removed (Stage~4);
\item the attention grid (Stage~5) and the patch grid.
\end{itemize}
Replies are compared as text, so ``identical'' means byte for byte.

\begin{table}[tbp]
\centering
\small
\begin{tabular}{lrrrrrrr}
\toprule
\textbf{tier} & \textbf{rows} & \textbf{reply} & \textbf{region} & \textbf{black} & \textbf{tokens} &
\textbf{attention} & \textbf{max $|\Delta|$} \\
\midrule
qwen7b-4bit & 40 & 40 & 40 & 40 & 40 & 40 & 0 \\
qwen7b-bf16 & 40 & 40 & 40 & 40 & 40 & 40 & 0 \\
\bottomrule
\end{tabular}
\caption{The shared backend against the frozen one on the A40 (job 488517, commit \texttt{89950a9},
\texttt{XAI\_REGION\_FRAME=rescaled}). Counts are rows identical in each field; max $|\Delta|$ is the
largest difference between the two attention grids. Rows cover all four rules (4-bit: 16, 7, 14, 3;
bf16: 20, 5, 10, 5), both verdicts, and about a third of rows whose black-masked reply differs from the
clean one, so the intervention paths were exercised, not only the clean one.
Files: \texttt{results/backend-parity/}.}
\label{tab:s29}
\end{table}

\textbf{Every field matched on every row, for both tiers} (Table~\ref{tab:s29}). The attention grids are
equal to the last bit. A smoke run of eight 4-bit rows on the RTX~3070 gave the same result before the
cluster run.

\textbf{One false start.} The first cluster job (488516) stopped at its GPU-test gate. The new GPU tests
for the shared backend run before the frozen backend's 8~GB budget test, and that test reads the
process-wide peak memory. It was measuring the previous test's model. The shared-backend tests now free
their model and reset the peak counter when they finish. No backend code changed between the two jobs.

\section{Standard vs.\ Non-Standard Choices}

\textbf{Standard.} A model-agnostic harness with per-model adapters is how evaluation suites such as
VLMEvalKit and lmms-eval are built.

\textbf{Non-standard, and stated in the paper.}
\begin{itemize}
\item The frozen backend is kept beside the shared one instead of being replaced, so the frozen results
are reproducible from the code that made them. The parity check is the evidence that the two agree.
\item Each model is asked for REGION in its own trained convention, not in one common unit. The question
and the rest of the format are identical.
\end{itemize}

\section{Implications}

\begin{itemize}
\item Adding a model is now writing one adapter of five duties. Nothing that computes a metric is touched.
\item The coordinate error of Chapter~\ref{chap:s2-8} cannot recur silently: every tier states its
convention, every convention has a round-trip test, and every new adapter is checked on rows drawn without
regard to their score before its first metric.
\item Next, queued on Scholar behind this check: the full battery for \texttt{qwen3b-bf16}; Stage~0 sizing
for \texttt{qwen32b-bf16} (two A40s) and \texttt{qwen32b-4bit}; the verdict-only battery for Llama~3.2
Vision; and the parity run of \texttt{qwen7b-bf16} in environment B (transformers 5.20.2, roster D5).
\end{itemize}
